\scanpage{330}\mspage{54}
% archive sheet number: 329

En effet, \deleted{cette} l'hypothèse \deleted{\unsure{sur}} $u\in W$ implique $\bar u\in W$, et l'hypothèse sur $\unsure{\mathcal{G}|G}$ implique dès lors $\bar u\in W_S^{\unsure{u|n}}$. Comme $W_S^u$ est un localiseur dans $M_{/S}$, on en conclut bien $u\in W_S^u$, qed.

Ainsi on a obtenu

\deleted{Nous pouvons appliquer ce lemme et}

\uline{Corollaire} 4.4.10. Supposons, \unsure{sous les} conditions \deleted{que} du lemme 4.4.9, que tout diagramme \inserted{comm.} triangulaire

\diagram{\begin{tikzcd}[ampersand replacement=\&,column sep=small]
 X\arrow[rr,dash,"{\unsure{u|du}}"]\arrow[dr,dash,"f"'] \& \& Y\arrow[dl,dash,"g"]\\
 \& S \&
\end{tikzcd}}

s'insère dans un diagramme \inserted{du type} (4.4.9.3) comme ci-dessus, (avec les relations explicitées sur les flèches qui y figurent). Alors
\[
 \deleted{\unsure{A}}B(\mathcal{G})\Longrightarrow \deleted{\unsure{A}}B(\unsure{\mathcal{F}|F}),
\]
Donc si de plus $\mathcal{G}\subset\mathcal{F}$, \deleted{\unsure{on a l}} de sorte que $A(\mathcal{F})\Longrightarrow \deleted{\unsure{B}}A(\mathcal{G})$ \deleted{est} tautologique, et de \={m} pour $B$ et pour $C$, on trouve
\[
 \begin{aligned}
 A(\mathcal{F})&\Longleftrightarrow A(\mathcal{G})\\
 B(\mathcal{F})&\Longleftrightarrow B(\mathcal{G})\\
 C(\mathcal{F})&\Longleftrightarrow C(\mathcal{G}).
 \end{aligned}
\]

Nous pouvons appliquer ce corollaire dans le cas qui nous occupe ici, avec
\[
 \deleted{\unsure{\mathcal{F}}}\mathcal{G}\overset{\mathrm{def}}{=}\mathcal{F}_{W_0}
 \subset\mathcal{F}\overset{\mathrm{def}}{\unsure{=|\,}}\mathrm{Fibf}_{W_0},
\]
